\documentclass[14pt]{extreport}
\usepackage[T1]{fontenc}
\usepackage[left=1in, right=1in, top=1in, bottom=1in]{geometry}
\usepackage{amsmath, amssymb, amsthm}
\usepackage{enumitem}
\usepackage{mathtools}
\usepackage{cancel}

% colored links
\usepackage{hyperref}
\hypersetup{
    colorlinks=true,
    linkcolor=blue,
    filecolor=magenta,      
    urlcolor=black!20!BurntOrange,
    }

% Inputting Python code
\usepackage[dvipsnames]{xcolor}
\definecolor{textblue}{rgb}{.2,.2,.7}
\definecolor{textred}{rgb}{0.54,0,0}
\definecolor{textgreen}{rgb}{0,0.43,0}
\usepackage{upquote}
\usepackage{listings}
\lstset{
    language=Python, 
    tabsize=4,
    basicstyle={\ttfamily},
    keywordstyle=\color{textblue},
    commentstyle=\color{textgreen},
    stringstyle=\color{textred},
    frame=none,
    columns=fullflexible,
    keepspaces=true,
    showstringspaces=false,
    xleftmargin=-15mm, % manual adjustment, figure out permanent solution
}

% Drawing figures
\usepackage{tikz}
\usetikzlibrary{matrix, positioning}
\usetikzlibrary{calc, shapes.symbols}

% Colored Boxes
\usepackage{tcolorbox}
\tcbuselibrary{skins,hooks}
\usetikzlibrary{shadows}
\usepackage{lipsum}

%Images
\usepackage{graphicx}
    \usepackage{subcaption}
    \usepackage{float}

%Formatting and Spacing
\setitemize[1]{noitemsep, parsep = 5pt, topsep = 5pt}
\setenumerate[1]{label = (\alph*), parsep = 1pt, topsep = 5pt}
\setlength\parindent{0pt}
\linespread{1.1}

%Cases Environment
\newlist{Cases}{enumerate}{3}
\setlist[Cases]{leftmargin = .25in, label = {Case \arabic*.}, topsep = 0.01in, itemsep = 0.04in, itemindent = .5in, parsep = 0in}

%Custom Title Fields
\newcommand{\lectTitle}{Lecture 11 Notes}
\newcommand{\lectTime}{March 7, 2022}
\newcommand{\lectClass}{Advanced Algorithms}
\newcommand{\lectClassInstructor}{Dana Randall}
\newcommand{\lectSection}{Spring 2024}
\newcommand{\lectAuthorName}{Sarthak Mohanty}


\title{
    \vspace{1.75in}
    \textbf{\lectClass:\\ \lectTitle}\\
    \vspace{0.1in}\large{\textit{\lectClassInstructor\ \lectSection}}
    \vspace{2in}
    \author{\textbf{Sarthak Mohanty}}
    \date{}
}

\begin{document}

\maketitle
\pagebreak

\section*{Ski-Rental Problem}
    Let's get some more practice with online algorithms. A person is going skiing for an unknown number of days $t$. Each day, they can either rent skis for $R$, or they can buy skis for $B$. If they rent skis, the person is presented with the choice the next day; if they buy, the game ends.

    \vspace{5mm}
    Since the game ends as soon as one buys, any algorithm for this problem can uniquely be defined with one number $k$ where $1 \le k \le t$, indicating the day where one buys.


    \vspace{5mm}
    Now suppose $k = B$. Then the total maximum cost is $k - 1 + B = 2B - 1$. In this cost, the offline cost is $\min(t, B) = B$. Thus, the competitive ratio is $(2B - 1) / B = 2 - \frac{1}{B}$.

    \vspace{5mm}
    Proof of optimality: Let $k$ be any strategy (buy after $k-1$ rents).
    \begin{itemize}
        \item Suppose you buy the skis at the $k$-th time and then break your leg and never ski again.
        \item Your total ski cost is $k-1+B$ and the optimum offline cost is $\min(k,B)$.
        \item For every $k$, the ratio $\frac{k-1+B}{\min(k,B)}$ is at least $(2-\frac{1}{B})$.
        \item Therefore, every strategy is at least $(2-\frac{1}{B})$-competitive.
    \end{itemize}

    
    The general rule: when balancing small incremental costs against a big one-time cost, you want to delay spending the big cost until you have accumulated roughly the same amount in small costs.

\subsection*{Random Algorithms}
    
\textbf{Theorem 1.1.} \textit{There exists a random algorithm that is $\left(\frac{e}{e-1}\right)$-competitive for the ski-rental problem.}

    \vspace{5mm}
    \textbf{Proof. (Sketch)} We assume there is a probability distribution $p_t$ associated with each day $t$, such that there is always a probability $p_t$ that you buy on day $t$. We must then calculate the theoretical value of $E[A(I)]$, $OPT(I)$ for every instance $I$. Then find the worst instance $I^*$ and optimize the probabilities $p$ for that instance.

    \vspace{5mm}
    We'll consider a simpler version of the problem where $b = 2$ and $l \leq 3$. The table below shows the competitiveness of each algorithm $A_i$ in each instance $I_j$, where $i$ determines which day to buy, and $j$ determines the last day of the ski tracks. Each entry of the table is the ratio between the algorithm and the optimal solution for that instance, i.e. $\frac{A_i(I_j)}{OPT(I_j)} \leq \alpha$.
    
    \begin{center}
    \begin{tabular}{c|cccc}
     & $A_1$ & $A_2$ & $A_3$ & $A_4$ \\
    \hline
    $I_1$ & 2 & 1 & 1 & 1 \\
    $I_2$ & 1 & $\frac{3}{2}$ & 1 & 1 \\
    $I_3$ & 1 & $\frac{3}{2}$ & 2 & $\frac{3}{2}$ \\
    \end{tabular}
    \end{center}
    
    The best deterministic algorithms are $A_2$ and $A_4$ that are both $\frac{3}{2}$-competitive. If we mix algorithms we could potentially improve the competitiveness. We can mix the two best algorithms. Consider that the algorithm $A$ runs algorithm $A_2$ and $A_4$, each with $\frac{1}{2}$ probability. Algorithm $A$ has the following results for the different instances:
    
    \begin{center}
    \begin{tabular}{c|c}
    $A$ &  \\
    \hline
    $I_1$ & 1 \\
    $I_2$ & $\frac{5}{4}$ \\
    $I_3$ & $\frac{3}{2}$ \\
    \end{tabular}
    \end{center}
    
    It is still $\frac{3}{2}$-competitive. We were unable to improve on the previous result. Note how algorithms $A_2$ and $A_4$ are simultaneously costly (in relation to the optimal) on instance $I_3$. Ideally, we'd choose two or more algorithms that are complementary to each other.

    \vspace{5mm}
    For this next try, we'll be choosing algorithms $A_1$ and $A_2$, note how they are complementary to each other. This time, we will not be attributing a priori probabilities to them. Let algorithm $A_1$ occur with probability $p$ and algorithm $A_2$ occur with probability $1 - p$. We have the following results for the algorithm $A$.
    
    \begin{center}
    \begin{tabular}{c|c}
    $A$ &  \\
    \hline
    $I_1$ & $p + 1$ \\
    $I_2$ & $\frac{3-p}{2}$ \\
    $I_3$ & $\frac{3-p}{2}$ \\
    \end{tabular}
    \end{center}
    
    To optimize the probability $p$ for our mixed strategy against a worst-case adversary, we apply the \textbf{Principle of Indifference}. We want to choose $p$ such that our algorithm's expected competitive ratio is identical across the relevant worst-case instances $I_1$ and $I_2$. By setting the expected ratios equal, $A(I_1) = A(I_2)$, we get:
    $$p + 1 = \frac{3-p}{2}$$
    Solving this yields $p = \frac{1}{3}$.
    
    \vspace{5mm}
    By taking the expected value of $A$, we obtain that the algorithm $A$ is $\frac{4}{3}$-competitive, a strict improvement on our previous findings.

    \vspace{5mm}
    \textbf{Main idea:} Yao's \textsc{Minimax} principle gives ``reduction'' to the deterministic case:

    \vspace{5mm}
    Randomized algorithms on worst-case deterministic instance $\approx$ Average-case complexity of deterministic algorithms

\vspace{5mm}
A randomized algorithm can be thought of as
\[
A = \left\{ \begin{array}{ccc}
A_1 & A_2 & \ldots \\
p_1 & p_2 & \ldots \\
\end{array} \right\},
\]

where \( A_i \) is a deterministic algorithm and \( A \) has probability \( p_i \) of running algorithm \( A_i \). We want to show that there is no randomized algorithm that is better than \( c^{-1} \)-competitive, e.g. there are no randomized algorithm that is 1.1-competitive, equivalently:

\vspace{5mm}
For every randomized algorithm $A$,  there exists $I$ such that  $$\mathbb{E}[A(I)] > 1.1 \times \text{OPT}(I)$$

\[
\Leftrightarrow \sum p_i A_i(I) > 1.1 \text{OPT}(I)
\]

A random instance \( I \) can be thought of as

\[
I = \left\{ \begin{array}{ccc}
I_1 & I_2 & \ldots \\
q_1 & q_2 & \ldots \\
\end{array} \right\},
\]

where \( I_j \) is a deterministic instance and \( I \) has probability \( q_i \) of instantiating to \( I_j \).

\vspace{5mm}
\textbf{Theorem 1.3.} (Yao's \textsc{Minimax} principle) For any randomized algorithm $A$ and random instance $I$ as detailed above we have

$$
\max_{I} \frac{\mathbb{E}[A(I)]}{\text{OPT}(I)} \geq \min_{A} \mathbb{E}_I \left[ \frac{A(I)}{\text{OPT}(I)} \right]
$$

Where on the left side we're running the randomized algorithm $A$ over the worst fixed instance and right side we're running the best deterministic algorithm over a random instance.

\vspace{5mm}
\textbf{Observation 1.} Under technical conditions the inequality becomes equality (non-trivial, uses, e.g., von Neumann's Minimax).

\vspace{5mm}
\textbf{Example:} Let's apply this to our matrix example. To prove a strong lower bound, we must act as the adversary and construct a worst-case probability distribution $I$ over the instances to maximize the \textit{minimum} expected ratio across all algorithms.

\vspace{5mm}
Notice that instance $I_2$ is strictly weaker than $I_3$ at punishing algorithms that buy late. Therefore, a smart adversary will ignore $I_2$ and focus on a mixture of the extreme cases: the short track ($I_1$) and the long track ($I_3$). Let the adversary play $I_1$ with probability $q$ and $I_3$ with probability $1-q$.

\vspace{5mm}
By the \textbf{Principle of Indifference}, we want to choose $q$ such that the algorithm designer is indifferent between their best ``buy early'' strategy ($A_1$) and their best ``buy late'' strategy ($A_4$). We set their expected ratios equal:
$$
\mathbb{E}[A_1] = \mathbb{E}[A_4] \implies q(2) + (1-q)(1) = q(1) + (1-q)\left(\frac{3}{2}\right)
$$
Solving for $q$:
$$
1 + q = 1.5 - 0.5q \implies 1.5q = 0.5 \implies q = \frac{1}{3}
$$

Thus, the optimal worst-case random instance is:
$$
I = \left\{ \begin{array}{cc}
I_1 & \text{with probability } \frac{1}{3} \\
I_2 & \text{with probability } 0 \\
I_3 & \text{with probability } \frac{2}{3} \\
\end{array} \right\}
$$

The expected competitive ratio is a weighted average over the possible instances. The table below shows the expected ratios of all deterministic algorithms against this specific random instance $I$:

\begin{center}
\begin{tabular}{c|cccc}
 & $A_1$ & $A_2$ & $A_3$ & $A_4$ \\
\hline
$\mathbb{E}_I \left[ \frac{A(I)}{\text{OPT}(I)} \right]$ & $\frac{4}{3}$ & $\frac{5}{3}$ & $\frac{5}{3}$ & $\frac{4}{3}$ \\
\end{tabular}
\end{center}

Against this adversary distribution, the absolute best a deterministic algorithm can do is a ratio of $4/3$ (achieved by choosing either $A_1$ or $A_4$). 

\vspace{5mm}
Because $\min_A \mathbb{E}_I \left[ \frac{A(I)}{\text{OPT}(I)} \right] = \frac{4}{3}$, Yao's Principle dictates that the right-hand side of the inequality evaluates to $4/3$. Therefore, we have rigorously proven that no randomized algorithm can have an expected competitive ratio better than $4/3$ for this scenario.
\end{document}